%-------------------------
% Cover letter -- ASML, AI | Computer Science internship: AI-driven software
% quality (Virtual Computing Platform, Veldhoven, J-00352656). Lives on this
% branch only.
%
% Same layout as the Philips letter: the resume's header, 1in sides, 0.75in
% top and bottom. Build: cd letters; pdflatex asml-cover-letter.tex
%
% Every project claim below is taken from cv/resume.tex on this branch.
% Style rules for this letter: plain words, no em-dashes, no "not X but Y".
%-------------------------

\documentclass[letterpaper,11pt]{article}

\usepackage[left=1in,right=1in,top=0.75in,bottom=0.75in]{geometry}
\usepackage[hidelinks]{hyperref}
\usepackage[english]{babel}

\urlstyle{same}
\pagestyle{empty}
\raggedright
\setlength{\parindent}{0pt}
\setlength{\parskip}{6pt}

% Same guard as resume.tex: machine-readable text under pdflatex, still
% compiles under XeTeX/tectonic.
\ifdefined\pdfgentounicode
  \input{glyphtounicode}
  \pdfgentounicode=1
\fi

\begin{document}

%-----------HEADING (identical to resume.tex)-----------
\begin{center}
    \textbf{\Huge \scshape Oleksii Krasnoshtanov} \\ \vspace{1pt}
    \small
    \href{mailto:oleksiikrasnoshtanov@gmail.com}{\underline{oleksiikrasnoshtanov@gmail.com}} $|$
    +31 6 29746862 $|$
    Tilburg, Netherlands \\ \vspace{2pt}
    \href{https://www.linkedin.com/in/oleksii-krasnoshtanov/}{\underline{linkedin.com/in/oleksii-krasnoshtanov}} $|$
    \href{https://github.com/alex-krasnoshtanov}{\underline{github.com/alex-krasnoshtanov}}
\end{center}
% The resume's section rule, reused as the header divider.
\vspace{-4pt}\noindent\rule{\textwidth}{0.4pt}

\today

ASML, Virtual Computing Platform team \\
Veldhoven, Netherlands

% Replace with the recruiter's or team lead's name if the portal shows one.
Dear Hiring Manager,

% A bare | in text prints as an em-dash under the default font encoding,
% hence \textbar{} in the job title.
I am applying for the AI \textbar{} Computer Science internship on AI-driven
software quality in the Virtual Computing Platform team (job ID J-00352656).
I am a third-year student in Applied Data Science and Artificial Intelligence
at Breda University of Applied Sciences.

I am highly passionate about ASML. Its work and research allow the world to
develop the most advanced technology we have, and I find the research behind
the company amazing. I want to grow my career in AI, and ASML is where I would
most like to do that.

I like that this assignment is practical. Working with software and test
engineers to find where AI can help, building prototypes and agents, and
measuring whether they make a real difference is the kind of work I enjoy most.
Software quality is also something I already pay attention to in my own
projects.

On a five-person platform project I owned deployment and CI/CD, with an 85\%
coverage gate, 430 tests, blue-green cutover and health-gated rollback. The
platform also tracked drift in confidence, latency and error rate, and
retrained the model based on researcher feedback on flagged predictions.
Finding problems early and feeding them back into the system is the part of
that project I would most like to build on.

% Agents are the core deliverable of this assignment, and nothing in the
% resume covers multi-step agents or tool use, so nothing is claimed here.
% Add a sentence only if you have real agent work to point to.
I also have hands-on experience with LLMs. In a two-person sprint I built a
self-hosted RAG assistant with Llama 3.3 70B, LangChain and FAISS, and the
client deployed it to production. It answers only from retrieved documents,
shows the sources behind each answer, and says clearly when the documents do
not cover a question. I think the same idea is important for AI that assesses
software issues: engineers need to see where a conclusion comes from before
they can trust it.

My current project gives me practice with the stakeholder side. As data
engineer on the Measurement Trust Project with CEVA Logistics and bol, I work
with several companies to get access to their data, help write business
requirements documents, and make sure our proposals are clear to non-technical
stakeholders.

% Availability: the posting wants a start as soon as possible, at least 5
% months, at least 4 days a week. Six months full-time covers the length and
% hours; the start date is the mismatch a reader will notice.
I work mainly in Python and use TypeScript for web frontends. I live in
Tilburg, so two days a week on site in Veldhoven is easy to arrange. I am
available full-time for six months from the beginning of 2027.

I would be happy to talk about how I could contribute to your team. Thank you
for your time and consideration.

Kind regards,

Oleksii Krasnoshtanov

\end{document}
